\poczatekwykladu{12}

\section{Całka a granica}
\label{sec:calka-a-granica}

Pracujemy na przestrzeni z miarą $(X,\mathcal F,\mu)$. Przypomnijmy,
że dla funkcji mierzalnych $f_n$ i $f$ zbieżność $f_n\to f$
$\mu$-prawie wszędzie oznacza
\[
  \mu(X\setminus E)=0,\qquad
  E=\left\{x\in X:\lim_{n\to\infty}f_n(x)=f(x)\right\}.
\]
Zbiór $E$ jest mierzalny. Mówimy wtedy, że $E$ jest zbiorem
\emph{miary pełnej}. Warunek ten implikuje $\mu(E)=\mu(X)$,
ale przy $\mu(X)=+\infty$ sama równość miar nie wystarcza
do stwierdzenia, że $E$ ma miarę pełną.

\begin{przyklad}[Granica punktowa a granica prawie wszędzie]
\label{prz:potegi-zbieznosc-prawie-wszedzie}
Na $[0,1]$ z miarą Lebesgue'a rozważmy $f_n(x)=x^n$. Wówczas
\[
  \lim_{n\to\infty}f_n(x)=f(x)=
  \begin{cases}
    0,&0\leq x<1,\\
    1,&x=1.
  \end{cases}
\]
Jeżeli $g\equiv0$ na $[0,1]$, to także $f_n\to g$
$\lambda$-prawie wszędzie, ponieważ $f$ i $g$ różnią się tylko
na singletonie $\{1\}$ miary zero.
\end{przyklad}

\section{Zbieżność monotoniczna: dowód}
\label{sec:dowod-zbieznosci-monotonicznej}

Wracamy do twierdzenia~\ref{tw:zbieznosc-monotoniczna}
Lebesgue'a o zbieżności monotonicznej (Beppo Levi). Dla mierzalnych
funkcji nieujemnych $f_n$ i $f$, przy $f_n\uparrow f$ prawie wszędzie,
zachodzi
\[
  \lim_{n\to\infty}\int_X f_n\,d\mu=\int_X f\,d\mu.
\]
Przed dowodem zobaczmy, dlaczego nie można pominąć warunku
monotoniczności ani zastąpić ciągu niemalejącego malejącym.

\begin{przyklad}[Przesuwający się przedział]
\label{prz:przesuwajacy-sie-przedzial}
Na $(\R,\mathcal L,\lambda)$ niech $f_n=\mathbf{1}_{[n,n+1]}$.
Dla każdego ustalonego $x\in\R$ mamy $f_n(x)=0$ od pewnego miejsca,
zatem $f_n\to0$ punktowo. Jednak
\[
  \int_\R\lim_{n\to\infty}f_n\,d\lambda=0,
  \qquad
  \lim_{n\to\infty}\int_\R f_n\,d\lambda=1.
\]
Przejście z granicą pod znak całki nie jest więc zawsze możliwe.
\end{przyklad}

\begin{przyklad}[Malejący ciąg na przestrzeni nieskończonej miary]
\label{prz:malejace-stale}
Na $(\R,\mathcal L,\lambda)$ niech $f_n(x)=1/n$ dla każdego $x$.
Mamy $f_n\downarrow0$ wszędzie, ale
\[
  \int_\R f_n\,d\lambda=+\infty\quad\text{dla każdego }n,
  \qquad
  \int_\R\lim_{n\to\infty}f_n\,d\lambda=0.
\]
Wniosek twierdzenia o zbieżności monotonicznej nie zachodzi zatem
dla dowolnego malejącego ciągu funkcji nieujemnych.
\end{przyklad}

Oba zjawiska przedstawia rysunek~\ref{fig:granica-calka-kontrprzyklady}.

\begin{figure}[htbp]
  \centering
  \includegraphics[width=0.98\textwidth]{granica-calka-kontrprzyklady.pdf}
  \caption{Przedział o stałym polu przesuwa się w prawo. W drugim przykładzie wysokości maleją do zera, lecz każda dodatnia stała ma na $\R$ całkę nieskończoną.}
  \label{fig:granica-calka-kontrprzyklady}
\end{figure}

\begin{proof}[Dowód twierdzenia~\ref{tw:zbieznosc-monotoniczna}]
Najpierw zakładamy $f_n\uparrow f$ w każdym punkcie $X$.
Z monotoniczności całki ciąg $\bigl(\int_X f_n\,d\mu\bigr)$ jest
niemalejący. Istnieje więc jego granica
\[
  L=\lim_{n\to\infty}\int_X f_n\,d\mu\in[0,+\infty].
\]
Ponieważ $f_n\leq f$, mamy $\int_X f_n\,d\mu\leq\int_X f\,d\mu$,
a po przejściu do granicy
\[
  L\leq\int_X f\,d\mu.
\]

Aby uzyskać przeciwną nierówność, korzystamy z definicji
\eqref{eq:calka-nieujemnej-supremum}. Wystarczy pokazać, że
$\int_X s\,d\mu\leq L$ dla każdej funkcji $s\in\mathrm{PMN}$
takiej, że $s\leq f$. Ustalmy taką funkcję $s$ oraz liczbę
$\alpha\in(0,1)$. Kładziemy
\begin{equation}
  B_n=\{x\in X:f_n(x)\geq\alpha s(x)\}.
  \label{eq:zbiory-progu-alpha}
\end{equation}
Zbiory $B_n$ są mierzalne, a z nierówności $f_n\leq f_{n+1}$
wynika $B_n\subseteq B_{n+1}$.

Pokażmy, że $\bigcup_{n=1}^\infty B_n=X$. Jeśli $s(x)=0$, to
$x\in B_n$ dla każdego $n$. Jeśli $s(x)>0$, to skończoność
wartości funkcji prostej i nierówność $s\leq f$ dają
\[
  \alpha s(x)<s(x)\leq f(x).
\]
Ponieważ $f_n(x)\uparrow f(x)$, od pewnego miejsca
$f_n(x)\geq\alpha s(x)$. Argument działa również dla $f(x)=+\infty$.
Schemat przekroczenia progu pokazano na rysunku~\ref{fig:prog-alpha-s}.

\begin{figure}[htbp]
  \centering
  \includegraphics[width=0.78\textwidth]{prog-alpha-s.pdf}
  \caption{Dla ustalonego $x$ z $s(x)>0$ wyrazy $f_n(x)$ od pewnego miejsca przekraczają próg $\alpha s(x)<f(x)$. Schemat dla skończonego $f(x)$.}
  \label{fig:prog-alpha-s}
\end{figure}

Dla dowolnego $A\in\mathcal F$ zachodzi $A\cap B_n\uparrow A$.
Ciągłość miary z dołu daje zatem
\begin{equation}
  \lim_{n\to\infty}\mu(A\cap B_n)=\mu(A).
  \label{eq:miara-przeciec-z-bn}
\end{equation}
Z definicji $B_n$ i nieujemności $f_n$ otrzymujemy
$f_n\geq\alpha s\mathbf{1}_{B_n}$ na całym $X$. Zapiszmy
$s=\sum_{i=1}^r a_i\mathbf{1}_{A_i}$, gdzie $a_i\geq0$,
a $A_i$ są mierzalne. Wtedy
\begin{align*}
  \alpha\sum_{i=1}^r a_i\mu(A_i\cap B_n)
  &=\alpha\int_X s\mathbf{1}_{B_n}\,d\mu\\
  &\leq\int_X f_n\mathbf{1}_{B_n}\,d\mu
  \leq\int_X f_n\,d\mu.
\end{align*}
Przechodząc z $n$ do nieskończoności i korzystając
z~\eqref{eq:miara-przeciec-z-bn}, dostajemy
\[
  \alpha\int_X s\,d\mu
  =\lim_{n\to\infty}\alpha\sum_{i=1}^r a_i\mu(A_i\cap B_n)
  \leq L.
\]
Suma jest skończona, a wszystkie składniki są nieujemne, więc
rachunek pozostaje poprawny również dla wartości $+\infty$.
Teraz pozwalamy $\alpha\uparrow1$:
\[
  \int_X s\,d\mu
  =\lim_{\alpha\uparrow1}\alpha\int_X s\,d\mu
  \leq L.
\]
Biorąc supremum po $s\in\mathrm{PMN}$, $s\leq f$, uzyskujemy
$\int_X f\,d\mu\leq L$. Wraz z pierwszą nierównością daje to
równość~\eqref{eq:zbieznosc-monotoniczna}.

Pozostaje wersja prawie wszędzie. Niech $E\in\mathcal F$ będzie
zbiorem miary pełnej, na którym $f_n(x)\uparrow f(x)$.
Definiujemy
\[
  \widetilde f_n=f_n\mathbf{1}_E,\qquad
  \widetilde f=f\mathbf{1}_E.
\]
Wówczas $\widetilde f_n\uparrow\widetilde f$ wszędzie: na $E$
wynika to z założenia, a poza $E$ wszystkie funkcje są zerowe.
Są one mierzalne i nieujemne. Ponadto
\[
  \int_X f\,d\mu=\int_X\widetilde f\,d\mu,
  \qquad
  \int_X f_n\,d\mu=\int_X\widetilde f_n\,d\mu,
\]
ponieważ całki po $X\setminus E$ są zerowe
(twierdzenie~\ref{tw:calka-po-zbiorze-miary-zero}). Stosując już
udowodnioną wersję punktową do $\widetilde f_n$, kończymy dowód.
\end{proof}

\section{Całkowanie szeregów funkcji nieujemnych}
\label{sec:calkowanie-szeregow-nieujemnych}

\begin{twierdzenie}[O całkowaniu szeregów]
\label{tw:calkowanie-szeregow-nieujemnych}
Niech $(f_n)$ będzie ciągiem nieujemnych funkcji mierzalnych na $X$.
Wtedy
\begin{equation}
  \int_X\sum_{n=1}^\infty f_n\,d\mu
  =\sum_{n=1}^\infty\int_X f_n\,d\mu.
  \label{eq:calkowanie-szeregow-nieujemnych}
\end{equation}
Obie strony mogą być równe $+\infty$.
\end{twierdzenie}

\begin{proof}
Sumy częściowe $g_n=\sum_{k=1}^n f_k$ są mierzalne i nieujemne.
Ponieważ $f_k\geq0$, mamy $g_n\leq g_{n+1}$ oraz punktowo
\[
  g_n\uparrow\sum_{k=1}^\infty f_k.
\]
Z definicji sumy szeregu o wyrazach nieujemnych, addytywności całki
dla skończonej sumy i twierdzenia o zbieżności monotonicznej wynika
\begin{align*}
  \sum_{n=1}^\infty\int_X f_n\,d\mu
  &=\lim_{n\to\infty}\sum_{k=1}^n\int_X f_k\,d\mu\\
  &=\lim_{n\to\infty}\int_X\sum_{k=1}^n f_k\,d\mu
  =\int_X\sum_{k=1}^\infty f_k\,d\mu.
\end{align*}
\end{proof}

\begin{wniosek}[Zamiana kolejności sumowania]
\label{wn:zamiana-kolejnosci-sum-nieujemnych}
Jeżeli $a_{m,n}\geq0$ dla $m,n\in\N$, to
\begin{equation}
  \sum_{m=1}^\infty\sum_{n=1}^\infty a_{m,n}
  =\sum_{n=1}^\infty\sum_{m=1}^\infty a_{m,n}.
  \label{eq:zamiana-kolejnosci-sum-nieujemnych}
\end{equation}
\end{wniosek}

\begin{proof}
Na $(\N,\mathcal P(\N),\eta)$, gdzie $\eta$ jest miarą liczącą,
kładziemy $f_n(m)=a_{m,n}$. Funkcje $f_n$ są mierzalne i nieujemne.
Całka względem miary liczącej jest sumą wartości funkcji
(przykład~\ref{prz:calka-miara-liczaca}). Stąd i z twierdzenia
\ref{tw:calkowanie-szeregow-nieujemnych} otrzymujemy
\begin{align*}
  \sum_{m=1}^\infty\sum_{n=1}^\infty a_{m,n}
  &=\int_\N\sum_{n=1}^\infty f_n\,d\eta\\
  &=\sum_{n=1}^\infty\int_\N f_n\,d\eta
  =\sum_{n=1}^\infty\sum_{m=1}^\infty a_{m,n}.
\end{align*}
\end{proof}

\section{Lemat Fatou}
\label{sec:lemat-fatou}

\begin{lemat}[Fatou]
\label{lem:fatou}
Niech $(f_n)$ będzie ciągiem nieujemnych funkcji mierzalnych na $X$.
Wtedy
\begin{equation}
  \int_X\liminf_{n\to\infty}f_n\,d\mu
  \leq\liminf_{n\to\infty}\int_X f_n\,d\mu.
  \label{eq:lemat-fatou}
\end{equation}
\end{lemat}

\begin{uwaga}[Nierówność ostra i granica górna]
\label{uw:fatou-nierownosc-ostra}
Nierówność w lemacie Fatou może być ostra. W przykładzie
\ref{prz:malejace-stale} dla $f_n\equiv1/n$ na $\R$ jej lewa strona
jest równa zero, a prawa $+\infty$.

Ten sam przykład wyklucza ogólne oszacowanie
\[
  \int_X\limsup_{n\to\infty}f_n\,d\mu
  \geq\limsup_{n\to\infty}\int_X f_n\,d\mu.
\]
Przeciwny kierunek również nie jest zawsze poprawny. Nawet
na przestrzeni skończonej miary nie musi zachodzić
\[
  \int_X\limsup_{n\to\infty}f_n\,d\mu
  \leq\limsup_{n\to\infty}\int_X f_n\,d\mu.
\]
Pokazuje to następny przykład.
\end{uwaga}

\begin{przyklad}[Wędrujące przedziały dyadyczne]
\label{prz:wedrujace-przedzialy-dyadyczne}
Rozważamy $[0,1]$ z miarą Lebesgue'a. Zaczynamy od
\begin{align*}
  f_1&=\mathbf{1}_{[0,1/2)},&
  f_2&=\mathbf{1}_{[1/2,1]},\\
  f_3&=\mathbf{1}_{[0,1/4)},&
  f_4&=\mathbf{1}_{[1/4,1/2)},\\
  f_5&=\mathbf{1}_{[1/2,3/4)},&
  f_6&=\mathbf{1}_{[3/4,1]}.
\end{align*}
Następnie przechodzimy kolejno przez osiem przedziałów długości
$1/8$, szesnaście przedziałów długości $1/16$ itd.
W każdym podziale ostatni przedział zawiera punkt $1$.

Dokładniej, dla $k\geq1$ i $j=0,\ldots,2^k-1$ definiujemy
\[
  I_{k,j}=
  \begin{cases}
    [j2^{-k},(j+1)2^{-k}),&0\leq j<2^k-1,\\
    [1-2^{-k},1],&j=2^k-1,
  \end{cases}
  \qquad f_{2^k-1+j}=\mathbf{1}_{I_{k,j}}.
\]
Początek ciągu przedstawia rysunek~\ref{fig:wedrujace-przedzialy-dyadyczne}.
Całki tworzą ciąg
\[
  \frac12,\ \frac12,\
  \underbrace{\frac14,\ \frac14,\ \frac14,\ \frac14}_{4\text{ wyrazy}},\
  \underbrace{\frac18,\ldots,\frac18}_{8\text{ wyrazów}},\ldots,
\]
więc $\int_{[0,1]}f_n\,d\lambda\to0$.
Każdy $x\in[0,1]$ należy do dokładnie jednego przedziału
w każdym podziale. Stąd $f_n(x)=1$ nieskończenie wiele razy,
a także $f_n(x)=0$ nieskończenie wiele razy. Zatem wszędzie
\[
  \limsup_{n\to\infty}f_n(x)=1,\qquad
  \liminf_{n\to\infty}f_n(x)=0.
\]
W szczególności
\[
  \int_{[0,1]}\limsup_{n\to\infty}f_n\,d\lambda=1
  >0=\limsup_{n\to\infty}\int_{[0,1]}f_n\,d\lambda.
\]
\end{przyklad}

\begin{figure}[htbp]
  \centering
  \includegraphics[width=0.76\textwidth]{wedrujace-przedzialy-dyadyczne.pdf}
  \caption{Przedziały, na których kolejne funkcje $f_n$ są równe $1$. Pełna kropka oznacza koniec należący do przedziału, pusta kropka koniec wyłączony.}
  \label{fig:wedrujace-przedzialy-dyadyczne}
\end{figure}

\begin{proof}[Dowód lematu~\ref{lem:fatou}]
Niech
\[
  g_n(x)=\inf\{f_n(x),f_{n+1}(x),\ldots\}
        =\inf_{k\geq n}f_k(x).
\]
Każda funkcja $g_n$ jest mierzalna jako infimum przeliczalnej rodziny
funkcji mierzalnych. Ponadto $g_n\geq0$, $g_n\leq g_{n+1}$,
$g_n\leq f_n$ i punktowo
\[
  \lim_{n\to\infty}g_n
  =\sup_{n\in\N}\inf_{k\geq n}f_k
  =\liminf_{n\to\infty}f_n.
\]
Z twierdzenia o zbieżności monotonicznej otrzymujemy
\[
  \lim_{n\to\infty}\int_X g_n\,d\mu
  =\int_X\lim_{n\to\infty}g_n\,d\mu
  =\int_X\liminf_{n\to\infty}f_n\,d\mu.
\]
Z drugiej strony $g_n\leq f_n$, więc monotoniczność całki daje
\[
  \liminf_{n\to\infty}\int_X f_n\,d\mu
  \geq\liminf_{n\to\infty}\int_X g_n\,d\mu
  =\lim_{n\to\infty}\int_X g_n\,d\mu.
\]
Po połączeniu obu wzorów otrzymujemy~\eqref{eq:lemat-fatou}.
\end{proof}
